
\begin{table}[htbp]
   
   \caption{\label{tab:fs_dwnstrm_minevio} First stage: effect of the Acid Rain Program on the number of upstream coal mines (averaged across upstream watersheds, utilities at most one HUC12 downstream)}
   \bigskip
   
   \centering
   
   \begin{adjustbox}{width = 0.45\textwidth, center}
      
      %start:tab
      
      \begin{tabular}{lr}
         \toprule
         
                                                  & num\_coal\_mines\_upstream\_mean\\      
         
                                                  & (1)\\  
         
         \midrule 
         post95 $\times$ sulfur\_unified\_mean    & -0.0822$^{***}$\\   
                                                  & (0.0222)\\   
         
         F-test (1st stage, clustered)            & 13.68\\  
         
          \\
         
         Observations                             & 6,225\\  
         R$^2$                                    & 0.73160\\  
         Within R$^2$                             & 0.02479\\  
         
         \bottomrule
         
      \end{tabular}
      
      %end:tab
      
      
   \end{adjustbox}
      {\tiny\linespread{1}\selectfont \par \raggedright 
   \textit{Notes:} The dependent variable is the number of coal mines in watersheds upstream of the utility's intake, averaged across upstream watersheds. The instrument interacts an indicator for the post-1995 period with the mean coal sulfur content of the intake watershed. The sample is utilities at most one watershed downstream of a coal mine. All specifications include utility and year fixed effects. Standard errors clustered at the utility level. Sample period 1985--2005. *** p$<$0.01, ** p$<$0.05, * p$<$0.1.}
   
\end{table}


